\todo{To PH: this section exceeds page budget. Please shrink to 2 pages and focu on 3 main points: 1) LLM in general, 2) LLMs for code, and 3) what is missing for HPC. }
The revolutionary breakthrough in machine learning and artificial intelligence significantly improves the maturity of natural language processing (NLP) by machine learning.
The potential and the impressive capabilities of machine learning were first revealed to the general public with the release of ChatGPT by OpenAI in 2022.  Large language models (LLMs) has become one of the most popular terms in research, industrial, news, and social media. 
\subsection{What is an LLM?}
Large language models (LLMs), a type of machine learning model trained by a large corpus of text data, have become the mainstay of NLP to perform tasks like text generation, summarization, and language translation.  
Various deep learning models, such as recurrent neural networks (RNNs) and Transformers, are developed to process sequential input data with target tasks text-to-text tasks.
Transformer models have become the primary deep-learning models for leading LLMs due to the adoption of a self-attention mechanism to enable higher parallelism in the model processing.
A typical Transformer architecture includes an encoder and a decoder. 
The input text of a Transformer model is first processed into a sequence of tokens by a tokenizer and further converted into an embedding represented by a vector.
The encoder maps the input embedding to a sequence of continuous representations (embedding). 
The decoder takes the output of the encoder together with the decoder output at the previous time step to generate an output sequence.
The output of the decoder is last processed to generate the prediction of the next word of the output sequence. 
\todo{A figure about large language model may be needed to help explain its components}
A Transformer-based LLM has the following major components:
\begin{itemize}
    \item Tokenization: A process, by a tokenizer, to convert the input text into a sequence of tokens that can be processed by the model.
    \item Embedding: Embeddings, represented by high-dimensional vector representation, are converted from the tokens to capture the semantic meaning of the input text.  
    \item Attention: The attention mechanism allows the model to assign different weights to tokens of the input sequence.  Therefore, the model can retrieve most relevant information and ignore less important details.
    \item Pre-training: Pre-training is the training process, both supervised or unsupervised, with a large-scale dataset. The goal of pre-training is to allow the models to capture the general patterns and the foundation of knowledge about the input texts or languages. 
    \item Transfer learning or fine-tuning: Transfer learning exploits a small but task-specific dataset to provide the model with more task-specific knowledge.
\end{itemize}


\subsection{Types of LLMs}
\todo{to PH: please shrink this subsection into one paragraph and merge into previous subsection}
With the design of the Transformer architecture, there are three types of LLM models with different designs and specific tasks:
\begin{itemize}
    \item Encoder models: Encoder models, or auto-encoding models, use only the encoder in a Transformer model, and the attention mechanism accesses all words of the input text sequence.  These models are most suitable to assist in the understanding of the text and specific tasks such as text classification and named entity recognition.  The most representative encoder models are the Bidirectional Encoder Representations from Transformers (BERT) and models derived from BERT, RoBERTa, ALBERT, and \textit{et al.}.
    \item Decoder models: Similarity, decoder models, or auto-regressive models use only the decoder from a Transformer model. However, the attention mechanism can access only words positioned before the word is processed in the current processing stage.  These models are most suitable for text-generation tasks.  The most representative decoder models are the models from the Generative Pre-trained Transformer (GPT) family, such as GPT, GPT2 \textit{et al.}.
    \item Encoder-decoder models: Encoder-decoder models, or sequence-to-sequence models, use both encoder and decoder from a Transformer model. There are attention layers from the encoder and decoder, respectively in an Encoder-decoder model.
    These models are best suitable for text generation tasks depending on a given input, such as summarization, translation and question answering.  T5 and BART are the most representative encoder-decoder models. 
\end{itemize}

\subsection{LLM examples}
The majority of the LLMs are general-purpose models trained to accomplish a wild range of tasks. As LLMs rely on a massive quantity of text data to provide knowledge and information regarding the language, semantics, and text contents. The datasets used for the training process greatly influence the skills and capabilities of the tasks executed by the LLMs.  Table~\ref{LLMs} summarizes representative LLMs, the training dataset used by the LLMs and the target tasks.


\begin{table}[htbp]
\centering
\caption{Example Language Models}
\label{modelExample}
\begin{tabular}{|l|l|r|l|r|c|}
\hline
\multicolumn{3}{|c}{\textbf{Model}} &
\multicolumn{2}{|c|}{\textbf{Training Data}} &
\multirow[t]{2}{*}{\textbf{Task}} \\
\cline{1-5}
\textbf{Type} & \textbf{Name} & \textbf{Size} & \textbf{Type} & \textbf{Size} &  \\ \hline
LM & BERT & 340M & G-NLT & 3.5B words & G-NLP \\ \hline
LM & CodeBERT & 125M & G-PLT & 2.1M  6.4M & G-PLP \\ \hline
LM & GraphCodeBERT &  & G-PLT &   & G-PLP \\ \hline
LLM & GPT-3 & 175B & G-NLT &  500B words  & G-NLP \\ \hline
LLM & LLaMA & 7\~65B & G-NLT &  500B words  & G-NLP \\ \hline
LLM & StarCoder & 15B & G-PLT &  1T tokens  & G-NLP \\ \hline
LLM & Dolly &  & &  & \\ \hline
\end{tabular}
\end{table}


\todo{PH: This table has too many empty space, not space efficient. Also need more models for coding tasks: e.g. CoDex,  }
\begin{table}[htbp]
\centering
\caption{Example Language Models}
\label{LLMs-old}
\begin{tabular}{|p{2.5cm}|p{3.5cm}|p{2.2cm}|p{3.5cm}|}
\hline
\textbf{Model/Size} & \textbf{Training Data} & \textbf{Data size} & \textbf{Tasks} \\ \hline
% we already have GPT-3, no need to list GPT-1
%GPT-1, 117M & 
%\begin{enumerate}
%    \item BookCorpus
%    \item 1B word benchmark 
%\end{enumerate}
%& ~2B words & 
%\begin{enumerate}
%    \item question answering
%    \item semantic similarity
%    \item classification
%\end{enumerate}
%\\ \hline
BERT, 340M & 
\begin{enumerate}
    \item BookCorpus
    \item English Wikipedia
\end{enumerate}  
& ~3.5B words &  diverse natural language understanding
tasks \\ \hline
GraphCodeBERT, 125M & 

& 2.1M bimodal /\newline 6.4M unimodal &  
\begin{enumerate}
    \item NL code search
    \item NL PL probing
    \item code documentation
\end{enumerate}
\\ \hline
CodeBERT, 125M & 
\begin{enumerate}
 \item codes in six programming languages \footnote{Python, Java,
JavaScript, PHP, Ruby, and Go}
\end{enumerate}
& 2.1M bimodal /\newline 6.4M unimodal &  
\begin{enumerate}
    \item NL code search
    \item NL PL probing
    \item code documentation
\end{enumerate}
\\ \hline
GPT-3, 175B & 
\begin{enumerate}
    \item Common Crawl
    \item WebText2
    \item Books1, Books2
    \item Wikipedia
\end{enumerate}
& ~500B words &  
\begin{enumerate}
    \item text prediction
    \item question answering
    \item translation
    \item reading comprehension
\end{enumerate}
\\ \hline
LLaMA, 7 to 65B & 
\begin{enumerate}
    \item CommonCrawl
    \item C4
    \item Github
    \item Books
    \item ArXiv
    \item StackExchange
\end{enumerate}
& 1.4T tokens & 
\begin{enumerate}
    \item common Sense Reasoning
    \item question answering
    \item reading comprehension
    \item mathematical reasoning
    \item code generation
\end{enumerate}
\\ \hline
StarCoder, 15B & The Stack 1.2 & 1 T tokens & 
\begin{enumerate}
    \item question answering
    \item tasks with code and text
\end{enumerate}
\\ \hline
\end{tabular}
\end{table}
\todo{merge two model tables into one, especially for large language models. The second table can keep information for smaller models we evaluated.}
Earlier generation LLMs were trained mainly by text data with a focus on natural language processing.  The sources of the training data mainly come from books, web content, and collected vocabulary from various natural languages. 
CodeBERT, derived from BERT, applied programming languages in training with the aim of bridging the gap between natural languages and programming languages.
Training data for the latest LLMs, such as LLaMA, started to include data representing programming languages and other artificial languages.  This also contributes to a significant growth in the size of the training data for modern LLMs.

LLMs are trained to have general text-related skills, such as text prediction/completion, reading comprehension, translation, and \textit{et al.}.  
LLMs that are capable for tasks related codes require training data associated with programming languages.  
Majority of the existing training data, for both natural languages and programming languages, are collected from public domains to provide general information and knowledge. 
Information and knowledge for a specific domain can be fed into the LLMs through the fine-tuning process. 
Without fine-tuning, LLMs might perform modestly for tasks that required specialized knowledge in a specific domain.  

To leverage the strength from LLMs for HPC research and developments, developing LLMs trained or fine-tuned with HPC-specific training data become a prerequisite.
However, a user would need to have sufficient knowledge of then given LLM, such as the pre-training data, fine-tuning data, and its capable skills, to effectively exploits the LLM for the tasks.
Without inspecting a LLM's detail and its training dataset, it is challenging to know a LLM's capability in any specific task. 
For example, a general LLM might not understand OpenMP and its language specification.  

Two thrusts can potentially improve LLMs' applicability of HPC workflows.  First, data curation for HPC training data will boost the LLMs' skills to cover more HPC-specific knowledge.  And second, simplifying and integrating end-to-end model inference for HPC workflow will greatly reduce the burden in HPC-specific developments.  With seamless end-to-end workflow, HPC users can easily apply and compare LLMs for a given specialized tasks.

Deploying Large Language Models (LLMs) for High-Performance Computing (HPC) tasks presents a unique set of challenges due to their specific requirements. \todo{What are these challenges?} Moreover, the rapid expansion of LLMs necessitates dedicated frameworks and APIs that can simplify model deployment and hasten inference implementation for the HPC community.
As a platform embraced by the Machine Learning (ML) community, \face offers effortless access to a wide variety of pre-trained models and datasets. However, its primary focus is on Natural Language Processing (NLP) limits its efficacy for HPC tasks.

\todo{background should mention what is missing for applying LLMs for HPC tasks (don't know which models, datasets, pipelines are available, lacking HPC specific data preprocessing, etc. Don't know which one is better than others for a certain task, etc.). Then we naturally transit to talk about our approach. This could be dedicated subsection.  }

\todo{1. LLM for general; 2. LLM for code; 3. LLM for HPC}
\todo{Hugging Face should be mentioned in background, then mention what is missing for HPC support. Currently, Hugging Face provides pipelines for NLP and image processing tasks only.  }
\todo{StarCode should be cited and mentioned in  text. Only listing it in the table without explanation is not sufficient}. 
